% Generated draft: requires mathematical and source review.
\documentclass[11pt]{article}
\usepackage[margin=25mm]{geometry}
\usepackage{fontspec}
\setmainfont{Latin Modern Roman}
\usepackage{amsmath,amssymb,mathtools,unicode-math}
\setmathfont{Latin Modern Math}
\usepackage{enumitem,needspace,xurl,hyperref}
\setlength{\parindent}{0pt}
\setlength{\parskip}{5pt}
\setlength{\emergencystretch}{3em}
\newcommand{\N}{\mathbb{N}}
\newcommand{\Z}{\mathbb{Z}}
\newcommand{\Q}{\mathbb{Q}}
\newcommand{\R}{\mathbb{R}}
\newcommand{\C}{\mathbb{C}}
\newcommand{\F}{\mathbb{F}}
\newcommand{\A}{\mathbb{A}}
\newcommand{\PP}{\mathbb{P}}
\newcommand{\E}{\mathbb{E}}
\newcommand{\eps}{\varepsilon}
\newcommand{\dd}{\,\mathrm{d}}
\begin{document}
\section*{UnsolvedMath ID: 700316}
\textbf{Source problem code:} catalogue entry
\subsection{Definitions and mathematical statement}
Let \(\rho,\tau\in\Z\), and let \(Q\) be the quaternion algebra over \(\Q\) with basis \(1,i,j,ij\), where \(i^2=\tau\), \(j^2=\rho\), and \(ij=-ji\). Its reduced norm is
\[
f(x,y,z,u)=x^2-\tau y^2-\rho z^2+\rho\tau u^2.
\]
Assume that \(f\) is indefinite over \(\R\) and anisotropic over \(\Q\); explicitly, it takes both positive and negative real values, and
\[
\forall(x,y,z,u)\in\Q^4,\qquad f(x,y,z,u)=0\Longrightarrow x=y=z=u=0.
\]
These assumptions imply \(\rho\tau\ne0\) and that \(\tau\) is not a square in \(\Q\). Put \(K=\Q(\sqrt{\tau})\), and use the representation
\[
X(x,y,z,u)=\begin{pmatrix}
x+\sqrt{\tau}\,y&\rho(z+\sqrt{\tau}\,u)\\
z-\sqrt{\tau}\,u&x-\sqrt{\tau}\,y
\end{pmatrix}\in M_2(K),
\qquad \det X=f(x,y,z,u).
\]
Let \(p\) be any prime with \(p\nmid\rho\tau\), and set
\[
R=\Z^{(p)}=\{m/p^t:m,t\in\Z\}=\Z[1/p],\qquad
\mathcal O=R\,1\oplus R\,i\oplus R\,j\oplus R\,ij.
\]
Define
\[
G=\{X(x,y,z,u):x,y,z,u\in R,\ \det X(x,y,z,u)=1\}.
\]
Here \(E\) denotes the identity matrix and \(Z(G)\) the center of \(G\).

The supplied conjecture is that, for every such choice of \(\rho,\tau,p\), every normal subgroup \(N\trianglelefteq G\) with \(N\not\subseteq Z(G)\) contains a full congruence subgroup
\[
N(\mathfrak A)=\{X\in G:X\equiv E\pmod{\mathfrak A}\}
\]
for some modulus \(\mathfrak A\). The source does not specify the ambient ring of \(\mathfrak A\), its admissible values, or the exact meaning of this matrix congruence, so this part of the formulation remains unresolved.

If the intended moduli are nonzero ideals \(I\) of \(R\), with congruence interpreted in the order \(\mathcal O\), a precise conditional formulation is obtained by defining
\[
\Gamma_{\mathcal O}(I)=\{X(x,y,z,u)\in G:x-1,y,z,u\in I\}
=\ker\bigl(G\longrightarrow(\mathcal O/I\mathcal O)^\times\bigr).
\]
Under that interpretation the conjecture reads
\[
\forall N\trianglelefteq G,\quad
N\not\subseteq Z(G)\Longrightarrow
\exists\,I\trianglelefteq R,\ I\ne(0),\quad
\Gamma_{\mathcal O}(I)\subseteq N.
\]
This interpretation of the moduli requires confirmation against the original source. Allowing the zero modulus would make the containment assertion vacuous. The source calls its stated normal-subgroup condition the congruence subgroup property and additionally asserts that this property is independent of the matrix representation of \(Q\); that assertion is retained as a source qualification.
\subsection{Short English statement}
For an indefinite quaternion norm that has no nonzero rational zero, form the group of norm-one elements whose four coordinates lie in \(\Z[1/p]\), where \(p\nmid\rho\tau\). The supplied conjecture asks whether every normal subgroup not contained in the center contains a full congruence subgroup. The source leaves the permitted congruence moduli unspecified; the displayed formulation using nonzero ideals of \(\Z[1/p]\) is conditional on that interpretation.
\subsection{Sources}
\begin{itemize}
\item Supplied local UnsolvedMath registry record: permanent ID \(700316\), exact source problem code \texttt{catalogue entry}, title \emph{Kourovka Notebook Problem 5.33}. The supplied record labels the problem \emph{open}; this is source metadata, not an independently verified status assessment. Upstream revision: \texttt{d371ee917ac9647fbb06cd6fc529deec6f3dda82}.
\item \emph{The Kourovka Notebook (Unsolved Problems in Group Theory)}, problem 5.33, fifth issue (1976); the supplied background identifies the twenty-first issue edition (2026), \texttt{arXiv:1401.0300v48}. The supplied collection description attributes the notebook's editing to E. I. Khukhro and V. D. Mazurov. This reference was supplied with the record and has not been independently consulted here.
\item Attribution is preserved in both supplied forms: the proposer metadata and background name J. Mennicke, while the mathematical statement begins with the attribution ``(Y. Ihara).'' Their respective roles are not explained in the supplied record.
\end{itemize}
\end{document}
